\documentclass[9pt,twocolumn,twoside]{pnas-new}

\templatetype{pnasresearcharticle}

\makeatletter
\setlength{\@fptop}{0pt}
\setlength{\@fpsep}{10pt}
\setlength{\@fpbot}{0pt plus 1fil}
\setlength{\@dblfptop}{0pt}
\setlength{\@dblfpsep}{10pt}
\setlength{\@dblfpbot}{0pt plus 1fil}
\makeatother

% Mark revised passages. To remove markings, redefine as: \newcommand{\new}[1]{#1}
\usepackage{xcolor}
\newcommand{\new}[1]{#1}

\begin{document}

\title{Regional Favoritism in Later Periods: Nightlights and Pollution in Leader Birth Regions}

\author[a,b]{Richard Schulz}
\author[a]{Muhammad Majid Altaf}
\author[a]{Ruby Wang}

\affil[a]{University of California, Berkeley, Berkeley, CA 94720, USA}
\affil[b]{ETH Zurich, 8092 Zurich, Switzerland}

\leadauthor{Schulz}

\significancestatement{Regional favoritism is a well-known finding in political economy: leaders tend to direct economic activity toward their birth regions. This paper studies how strongly the canonical birth-region nightlights result appears in the later annual windows used for our environmental extension. In those later windows, the reduced form is much weaker than in the long DMSP replication window. The attenuation is broad across geographic, political, and measurement-motivated subsamples. Pollution outcomes do not suggest that leaders shifted from directing more growth to directing cleaner growth toward birth regions.}

\authorcontributions{R.S. conceived the study, built and integrated the main data pipeline, ran the core empirical analysis, and drafted the manuscript. M.M.A. contributed to the NO$_2$ analysis. R.W. contributed to manuscript review and proofreading.}
\authordeclaration{The authors declare no competing interests.}
\correspondingauthor{To whom correspondence should be addressed. E-mail: rschul@ethz.ch}
\keywords{regional favoritism $|$ political economy $|$ nighttime lights $|$ air pollution $|$ autocracy}
\acknow{We thank Professor Joshua Blumenstock for feedback. All errors are our own.}

\begin{abstract}
\new{The canonical finding that political leaders direct economic activity toward their birth regions, identified in nightlights, is much weaker in post-2005 annual windows than in the long DMSP replication sample. Using a global ADM2 district-year panel that combines PLAD leader birthplaces, harmonized nightlights, and surface pollution measures, with ADM2 and country-year fixed effects, we recover the classic result in 1992--2013 (lag-0 coefficient 0.0128, $p = 0.030$) but estimate near-zero and statistically uninformative effects in DMSP 2005--2013 (0.0017, $p = 0.768$) and harmonized 2005--2019 (0.0004, $p = 0.966$). The attenuation survives targeted cuts by geography, leader tenure, autocracy, luminosity headroom, and pollution-support samples; Africa and Sub-Saharan Africa drive the historical effect but no post-2005 subsample restores it. We then test whether favoritism shifted from the quantity to the composition of growth, a green favoritism hypothesis under which birth regions would receive cleaner activity even when nightlights do not rise. The pollution data do not support this interpretation: pooled NO$_2$ and pollution-intensity estimates are zero to positive, and PM$_{2.5}$ results are also not negative. Our contribution is diagnostic: we document that post-2005 attenuation is broad rather than an artifact of a single endpoint, and we rule out a simple shift toward greener birth-region growth as its explanation.}
\end{abstract}

\dates{Compiled on \today}
\doi{\url{www.pnas.org/cgi/doi/10.1073/pnas.XXXXXXXXXX}}

\maketitle
\thispagestyle{firststyle}
\ifthenelse{\boolean{shortarticle}}{\ifthenelse{\boolean{singlecolumn}}{\abscontentformatted}{\abscontent}}{}

\new{Political leaders direct visible economic activity toward their birth regions, a pattern documented across countries using subnational nightlights \cite{HodlerRaschky2014}. Whether that pattern persists in the more recent annual windows used by environmental and harmonized lights data, roughly the post 2005 era, is a separate question. These windows matter because they overlap with global pollution measures and with the harmonized DMSP to VIIRS panel that successor work in this literature increasingly relies on \cite{Bora2025, DubenHodlerRaschky2026}. They are also exactly where the canonical reduced form becomes much harder to detect.}

\new{This paper makes three contributions. First, we systematically document the post 2005 attenuation of the birth region nightlights effect across a window grid and across theory favored subsamples, including geography, leader tenure, autocracy, luminosity headroom, and pollution support countries, and show that no cut restores a precise activity effect. Second, we test whether the attenuation reflects a shift from the quantity to the composition of favored growth. Under a green favoritism hypothesis, leaders direct cleaner rather than greater economic activity to their birth regions, and pollution intensity should fall in birth regions even when nightlights do not rise.This possibility matters for environmental inequality as well as distributive politics \cite{CopelandTaylor2004}, but NO$_2$ and PM$_{2.5}$ outcomes do not support it. Third, we situate the findings against recent replication work using extended pooled samples \cite{Bora2025, DubenHodlerRaschky2026}: our own window-grid evidence (Figure~\ref{fig:windowgrid}) shows that positive coefficients concentrate in earlier DMSP windows, raising the possibility that long pooled samples spanning both early and post-2005 years can mask the period-specific attenuation we document.}

\new{The lag-0 birth-region coefficient is 0.0128 in DMSP 1992--2013 ($p = 0.030$) and remains positive in nearby historical windows. In DMSP 2005--2013 it falls to 0.0017 ($p = 0.768$), and in harmonized 2005--2019 it is 0.0004 ($p = 0.966$). Africa and Sub-Saharan Africa carry the long-window effect, with coefficients near 0.044, but these subsamples do not rescue the post-2005 estimate. Pooled post-2005 NO$_2$ and pollution-intensity coefficients are positive rather than negative. We read the pollution evidence as suggestive against green favoritism rather than as a standalone positive finding, since precision falls under country-level clustering. The historical 95\% confidence intervals still contain the post-2005 point estimates, so the result is a loss of detectable signal, not proof of zero.}

\section{Data and Empirical Design}

The analysis combines five core inputs. First, the Political Leaders' Affiliation Database (PLAD) provides leader identity, spell timing, and subnational birthplace information used to code treatment \cite{Bomprezzi2025PLAD}. Second, nightlights provide the economic-activity proxy. We use a replication-style DMSP ADM2 panel for 1992--2013 and a harmonized DMSP-to-VIIRS panel for 2005--2019, with DMSP-based observations through 2013 and simulated VIIRS afterward \cite{Elvidge2013VIIRS}. DMSP records luminance as an integer Digital Number (DN) from 0 to 63, with known saturation near the ceiling \cite{Gibson2021}. Third, we use ACAG surface NO$_2$ data for the 2005--2019 pollution window \cite{Cooper2022NO2}. Fourth, we use ACAG SatPM PM$_{2.5}$ as a longer-window pollution extension. Fifth, GADM ADM2 boundaries provide the common subnational geography \cite{GADM2022}.

Treatment is defined at the district-year level. A district is treated when it is the birth region of the national leader in office in that year. To ensure a unique treatment assignment in transition years, we truncate the earlier of any two consecutive leader spells that overlap within a country.

The core estimating equation uses ADM2 fixed effects and country-year fixed effects:
\[
\log(Y_{i c t} + 0.01) = \beta \, BirthRegionLeader_{i c t} + \alpha_i + \gamma_{c t} + \varepsilon_{i c t},
\]
where $i$ indexes ADM2 regions, $c$ countries, and $t$ years. With ADM2 and country-year fixed effects, $\beta$ is identified by changes in the current leader's birth district relative to other districts in the same country-year, net of permanent district differences. A causal interpretation requires birth and non-birth districts to have followed parallel within-country trends absent leader ascension. That assumption could fail if future leaders come from districts already on different trajectories or if national transitions coincide with region-specific shocks, so we interpret the estimates as reduced-form evidence rather than a fully causal decomposition.

Standard errors are clustered by leader spell in the baseline, following the leader-term clustering convention in recent replication work such as \cite{Bora2025}. Because \cite{HodlerRaschky2014} and \cite{DubenHodlerRaschky2026} use country-level clustering, we also report country-clustered robustness checks for the main estimates. Following the baseline specification in \cite{HodlerRaschky2014}, we estimate lag 0, lag 1, and lag 2 treatment indicators, and focus on lag 0 in the main text as the most direct test of the regional favoritism hypothesis.

Our later-period environmental outcomes are log nightlights, log NO$_2$, and a pollution-intensity measure defined as $\log(\text{NO}_2 + 0.01) - \log(\text{Nightlights} + 0.01)$. Because both component regressions use the same regressor and fixed effects, the intensity coefficient is mechanically a linear combination of the NO$_2$ and nightlights coefficients ($\beta_{intensity} = \beta_{NO2} - \beta_{lights}$). We report this measure as a compact summary of whether pollution rises more or less than nightlights in birth regions. Because it is constructed directly from the pollution and nightlights outcomes, it is not an independent third test. In the PM$_{2.5}$ extension, the analogous outcome is $\log(\text{PM}_{2.5} + 0.01) - \log(\text{Nightlights} + 0.01)$.

\section{Nightlights Reduced Form and Attenuation}

\subsection*{Replication in the long DMSP era}

The long DMSP replication remains intact. In the replication-style 1992--2013 sample, the lag-0 birth-region coefficient is 0.0128 with $p = 0.030$ (Table~\ref{tab:firststage}). This corresponds to a roughly 1.3\% increase in nightlights in treated birth regions, in line with magnitudes in the original regional-favoritism literature \cite{HodlerRaschky2014, DeLuca2018}. The same coefficient is larger in 1995--2013 (0.0187, $p = 0.001$) and remains positive, though weaker, in 2000--2013 (0.0110, $p = 0.056$). The baseline activity effect is therefore not a narrow late-1990s artifact.

\subsection*{Attenuation in later windows}

The key change is that the nightlights reduced-form point estimate is much smaller and statistically uninformative in later windows. In 2005--2013, still within the DMSP era, the lag-0 coefficient is 0.0017 with $p = 0.768$. In the harmonized 2005--2019 sample, it is 0.0004 with $p = 0.966$. The confidence intervals still contain the historical estimate, so the evidence should be read as a loss of detectable signal rather than proof that the effect is exactly zero.

The window-grid diagnostics in Figure~\ref{fig:windowgrid} reinforce that reading. Positive coefficients are concentrated in DMSP-only windows with earlier starts and end years around the late 2000s or early 2010s. Harmonized post-2005 windows, by contrast, remain small and statistically uninformative throughout. The reduced-form problem is not a single bad endpoint choice. It reflects a broader pattern of later-period attenuation.

\subsection*{Additional diagnostics on later-period attenuation}

Additional diagnostics narrow the plausible explanations for that attenuation. First, the null is not obviously a pure harmonization artifact. A VIIRS-only check on the standalone 2018--2023 panel is also near zero: the lag-0 coefficient is 0.0038 with $p = 0.722$. Second, the political composition of the sample does not shift enough to make a simple ``later years are much more democratic'' story convincing. The share of leader-years in the strict-autocracy bin falls only modestly between the long and later windows. The more visible change is shorter observed spell support: average active years per spell fall from 5.22 in 1992--2013 to 3.97 in 2005--2013. This spell-length comparison is partly mechanical because the narrower 2005--2013 window truncates observed active years for spells starting early or ending late.

Third, treated districts in the later DMSP years are already much brighter than untreated districts. In 2005--2013, the median treated district has DMSP luminosity of 16.84, compared with 8.83 for untreated districts, and the treated upper tail sits near the DMSP ceiling at about 61. That is consistent with mechanical attenuation in later years if leaders are increasingly born in already-bright districts. In line with this, the average ascension-year within-country light rank of birth districts is high in both periods and slightly higher in the later sample. Treatment geocoding quality, on the other hand, is stable across periods, so worsening PLAD birthplace measurement is not the leading explanation.

\begin{table*}[t]
\centering
\caption{Lag-0 nightlights activity-effect comparison across key windows}
\begin{tabular*}{\textwidth}{l @{\extracolsep{\fill}} r r r r r r}
Window & Coef. & SE & $p$-value & N & Treated obs. & Clusters \\
\midrule
1992--2013 & 0.0128 & 0.0059 & 0.030 & 951{,}915 & 2{,}522 & 2{,}723 \\
1995--2013 & 0.0187 & 0.0055 & 0.001 & 829{,}988 & 2{,}196 & 2{,}358 \\
2000--2013 & 0.0110 & 0.0058 & 0.056 & 615{,}660 & 1{,}647 & 1{,}759 \\
2005--2013 & 0.0017 & 0.0058 & 0.768 & 398{,}636 & 1{,}068 & 1{,}186 \\
2005--2019 & 0.0004 & 0.0087 & 0.966 & 671{,}399 & 1{,}701 & 1{,}923 \\
\bottomrule
\end{tabular*}
\addtabletext{All specifications use ADM2 fixed effects, country-year fixed effects, and leader-spell clustering. The 2005--2019 row uses the harmonized DMSP--VIIRS nightlights panel.}
\label{tab:firststage}
\end{table*}

\begin{table}[t]
\centering
\caption{Robustness to country-level clustering}
\scriptsize
\begin{tabular*}{\columnwidth}{@{\extracolsep{\fill}}lrrrr}
Spec. & Coef. & Lead. SE & Ctry. SE & Ctry. $p$ \\
\midrule
NL, 1992--2013 & 0.0128 & 0.0059 & 0.0075 & 0.087 \\
NL, 2005--2013 & 0.0017 & 0.0058 & 0.0081 & 0.832 \\
NL, 2005--2019 & 0.0004 & 0.0087 & 0.0107 & 0.973 \\
NO$_2$, 2005--2019 & 0.0308 & 0.0154 & 0.0204 & 0.132 \\
NO$_2$ int., 2005--2019 & 0.0333 & 0.0179 & 0.0239 & 0.163 \\
\bottomrule
\end{tabular*}
\addtabletext{NL denotes nightlights. All rows use ADM2 fixed effects and country-year fixed effects. Baseline standard errors cluster by leader spell; the robustness column clusters by country, following the original Hodler--Raschky convention.}
\label{tab:cluster-robustness}
\end{table}

\begin{figure*}[t]
\centering
\includegraphics[width=\textwidth]{figures/first_stage_window_grid.png}
\caption{Lag-0 birth-region nightlights coefficients across valid start and end windows. Each cell is a separate ADM2 fixed-effect and country-year fixed-effect regression; color encodes the birth-region coefficient. Earlier DMSP windows show a block of positive estimates, while post-2005 harmonized windows remain near zero.}
\label{fig:windowgrid}
\end{figure*}

\section{Heterogeneity in Attenuation}

If the pooled post-2005 activity reduced form masks heterogeneity, the most plausible places to look are settings where the original favoritism mechanism should be strongest or easiest to measure. We therefore test four families of subsamples: geography, political capacity, measurement headroom, and pollution-support samples. The geography cut captures Hodler and Raschky's finding of especially strong effects in Sub-Saharan Africa. The political-capacity cut isolates long-tenure leaders and long-tenure autocrats. The headroom splits ask whether DMSP saturation hides effects in already-bright districts. The pollution-support splits ask whether environmental outcomes are only informative in countries with enough NO$_2$ variation.

This exercise yields one historical result worth noting but no post-2005 rescue. In the full 1992--2013 DMSP window, the Africa coefficient is 0.0428 ($p = 0.009$) and the Sub-Saharan Africa coefficient is 0.0442 ($p = 0.011$), while the non-SSA coefficient is small and insignificant. The long-run replication is thus concentrated in the geography emphasized by the original literature. In 2005--2013, however, the Africa and SSA estimates turn negative and insignificant. The other theory-favored post-2005 samples in Table~\ref{tab:subsample-firststage} also fail to revive the activity reduced form: long-tenure autocrats are consistent with zero and too imprecise to distinguish from the historical effect, dropping high-luminosity districts does not help, and the high-NO$_2$ country sample remains near zero.

The harmonized 2005--2019 panel gives the most suggestive, though still inconclusive, later pattern. Africa (0.0262, $p = 0.167$) and Sub-Saharan Africa (0.0315, $p = 0.135$) are larger than the pooled null, but neither clears our minimum-support rule of a positive coefficient above 0.005 with $p < 0.10$. We therefore treat pollution outcomes in these samples as diagnostic rather than confirmatory (Table~\ref{tab:subsample-pollution}).

\begin{table*}[t]
\centering
\caption{Targeted activity-effect subsample search}
\begin{tabular*}{\textwidth}{l @{\extracolsep{\fill}} l r r r r r}
Sample & Window & Coef. & SE & $p$-value & N & Treated \\
\midrule
Africa & DMSP 1992--2013 & 0.0428 & 0.0163 & 0.009 & 109{,}126 & 705 \\
Sub-Saharan Africa & DMSP 1992--2013 & 0.0442 & 0.0174 & 0.011 & 63{,}378 & 631 \\
Non-SSA & DMSP 1992--2013 & 0.0053 & 0.0059 & 0.370 & 888{,}537 & 1{,}891 \\
\midrule
Pooled & DMSP 2005--2013 & 0.0017 & 0.0058 & 0.768 & 398{,}636 & 1{,}068 \\
Africa & DMSP 2005--2013 & -0.0052 & 0.0143 & 0.714 & 46{,}939 & 306 \\
Sub-Saharan Africa & DMSP 2005--2013 & -0.0063 & 0.0164 & 0.700 & 28{,}191 & 279 \\
Long-tenure leaders & DMSP 2005--2013 & -0.0062 & 0.0087 & 0.479 & 300{,}372 & 862 \\
Long-tenure autocrats & DMSP 2005--2013 & 0.0081 & 0.0173 & 0.637 & 42{,}717 & 261 \\
Drop top 10\% luminosity & DMSP 2005--2013 & 0.0010 & 0.0082 & 0.899 & 357{,}745 & 613 \\
Leader birth DN $< 61$ & DMSP 2005--2013 & 0.0020 & 0.0063 & 0.750 & 373{,}862 & 1{,}034 \\
Low birth brightness & DMSP 2005--2013 & 0.0016 & 0.0108 & 0.880 & 142{,}320 & 634 \\
High baseline NO$_2$ countries & DMSP 2005--2013 & -0.0012 & 0.0069 & 0.866 & 312{,}633 & 603 \\
\midrule
Africa & Harmonized 2005--2019 & 0.0262 & 0.0189 & 0.167 & 83{,}101 & 511 \\
Sub-Saharan Africa & Harmonized 2005--2019 & 0.0315 & 0.0210 & 0.135 & 52{,}001 & 466 \\
Low NO$_2$ dispersion countries & Harmonized 2005--2019 & 0.0219 & 0.0183 & 0.232 & 95{,}961 & 621 \\
\bottomrule
\end{tabular*}
\addtabletext{All rows use lag-0 treatment, ADM2 fixed effects, country-year fixed effects, and leader-spell clustering. Long tenure means total PLAD spell length of at least five years. Long-tenure autocracy additionally requires $v2x\_polyarchy < 0.3$. The minimum-support rule is coef. $\geq 0.005$, $p < 0.10$, and at least 50 treated observations in a post-2005 sample. No post-2005 subsample satisfies that rule.}
\label{tab:subsample-firststage}
\end{table*}

\begin{figure*}[t]
\centering
\includegraphics[width=0.74\textwidth]{figures/final_subsample_first_stage_plot.png}
\caption{Targeted post-2005 activity-effect estimates. Points are lag-0 coefficients and horizontal bars are 95\% confidence intervals for pooled, Africa, Sub-Saharan Africa, long-tenure, long-tenure-autocracy, luminosity-headroom, low-brightness, and high-NO$_2$ country samples. The pooled harmonized 2005--2019 estimate is included as a reference row.}
\label{fig:subsample-firststage}
\end{figure*}

\begin{table*}[t]
\centering
\caption{Pollution outcomes in strongest positive diagnostic samples}
\begin{tabular*}{\textwidth}{l @{\extracolsep{\fill}} r r r r r r}
Sample & Activity coef. & Activity $p$ & NO$_2$ intensity & NO$_2$ $p$ & PM$_{2.5}$ intensity & PM$_{2.5}$ $p$ \\
\midrule
Sub-Saharan Africa & 0.0315 & 0.135 & -0.0323 & 0.237 & -0.0336 & 0.127 \\
Africa & 0.0262 & 0.167 & -0.0267 & 0.295 & -0.0267 & 0.178 \\
Low NO$_2$ dispersion & 0.0219 & 0.232 & 0.0156 & 0.685 & -0.0203 & 0.281 \\
\bottomrule
\end{tabular*}
\addtabletext{The activity columns use the harmonized 2005--2019 nightlights panel. Pollution-intensity outcomes are log pollution minus log nightlights in the same subsample. These are diagnostic estimates because none of the activity effects is statistically significant under the stated criterion.}
\label{tab:subsample-pollution}
\end{table*}

\section{Pollution-Side Results}

\subsection*{Pollution outcomes in the 2005--2019 sample}

Pollution estimates test whether the post-2005 attenuation reflects a shift from the quantity to the composition of favored growth. In the pooled 2005--2019 sample (Table~\ref{tab:pollution}), the lag-0 nightlights coefficient is near zero, the lag-0 NO$_2$ coefficient is positive at 0.0308 ($p = 0.045$), and the lag-0 pollution-intensity outcome is also positive at 0.0333 ($p = 0.063$). These point estimates go against a cleaner-growth interpretation. At the same time, Table~\ref{tab:cluster-robustness} shows that the NO$_2$ estimates become less precise under country clustering, so they are better read as suggestive evidence against green favoritism than as a standalone positive pollution finding.

\begin{table*}[t]
\centering
\caption{Main later-period pollution results}
\begin{tabular*}{\textwidth}{l @{\extracolsep{\fill}} l r r r r}
Outcome & Window / spec & Coef. & SE & $p$-value & N \\
\midrule
$\log(\text{Nightlights}+0.01)$ & 2005--2019 lag 0 & -0.0026 & 0.0091 & 0.778 & 631{,}377 \\
$\log(\text{NO}_2+0.01)$ & 2005--2019 lag 0 & 0.0308 & 0.0154 & 0.045 & 631{,}377 \\
NO$_2$ pollution intensity & 2005--2019 lag 0 & 0.0333 & 0.0179 & 0.063 & 631{,}377 \\
PM$_{2.5}$ pollution intensity & 2005--2019 lag 0 & 0.0041 & 0.0095 & 0.665 & 669{,}438 \\
PM$_{2.5}$ pollution intensity & 2005--2019 lag 1 & 0.0122 & 0.0078 & 0.120 & 669{,}438 \\
PM$_{2.5}$ pollution intensity & 2005--2019 lag 2 & 0.0175 & 0.0074 & 0.019 & 669{,}438 \\
\bottomrule
\end{tabular*}
\addtabletext{Pollution intensity is $\log(Pollution+0.01) - \log(Nightlights+0.01)$. Negative coefficients would indicate cleaner growth in leader birth regions; current estimates are zero to positive.}
\label{tab:pollution}
\end{table*}

\subsection*{Autocracy and PM$_{2.5}$ extensions}

The heterogeneity checks do not recover the later activity effect either. Restricting the sample to country-years with $v2x\_polyarchy < 0.3$ yields a lag-0 nightlights coefficient of 0.0111 ($p = 0.574$) in the DMSP 1992--2013 sample and 0.0076 ($p = 0.552$) in the DMSP 2005--2013 sample. These estimates are positive but too imprecise to rescue the historical result.

The PM$_{2.5}$ extension tells the same story. In longer windows, PM$_{2.5}$ level regressions are positive rather than negative, and in the overlapping 2005--2019 sample the PM$_{2.5}$ pollution-intensity outcome is also not negative. The lag-2 estimate in Table~\ref{tab:pollution} is positive and statistically significant, though we read it cautiously given multiple testing across lags. On balance, pollution outcomes do not explain the weaker post-2005 nightlights result as a shift toward cleaner birth-region growth.

\section{Discussion}

The main result is that the canonical birth-region nightlights reduced form is much weaker in later periods, especially after 2005. The long DMSP replication still holds, especially in Africa and Sub-Saharan Africa, but the later-period estimates are close to zero and statistically uninformative in the pooled sample and in the most natural diagnostic subsamples. This points to a change over time in the strength of the original finding, not just to a change in how nightlights are measured.

The attenuation is broad enough that it cannot be dismissed as a single endpoint choice. We find suggestive patterns in shorter observed leader spells and reduced luminosity headroom, but those explanations are not definitive: the subsample tests dropping high-luminosity districts do not recover the effect, and spell-length attenuation is partly mechanical in the shorter window. The 95\% confidence intervals for the post-2005 coefficients in Table~\ref{tab:firststage} still contain the historical 1.3\% effect, so we cannot statistically rule out a constant-effect null. The results indicate a loss of detectable signal, not its proven absence.

The pollution analysis narrows one explanation for that attenuation. If post-2005 favoritism shifted from more growth to cleaner growth, pollution intensity should fall in leader birth regions even when nightlights do not rise. Instead, the pooled pollution estimates are zero to positive, and the negative Africa and Sub-Saharan Africa intensity point estimates in Table~\ref{tab:subsample-pollution} are too imprecise to carry the cleaner-growth interpretation. The paper therefore does not solve why the reduced form weakens after 2005, but it shows that the weakening is broad and is not explained by a simple shift toward greener birth-region favoritism.

\showacknow

\section*{Data, Materials, and Code Availability}

Code, notebooks, and manuscript sources are available in the public GitHub repository at \url{https://github.com/RichSchulz/political-leader-pollution}. The reported tables are produced from the analysis notebooks in \texttt{analysis/}, including \texttt{replication.ipynb}, \texttt{green\_favoritism.ipynb}, \texttt{green\_favoritism\_pm25.ipynb}, \texttt{final\_subsample\_search.ipynb}, and \texttt{final\_clustering\_robustness.ipynb}. External inputs include PLAD, GADM, ACAG, V-Dem, and harmonized nightlights data obtained under the providers' access conditions.

\bibliography{references}

\end{document}
